% !TeX root = ../main.tex

\chapter{绪论}

\section{研究背景}

\subsection{新时代信息化人才培养的需要}

2018 年《教育信息化 2.0 行动计划》明确指出，人才培养需实现从技术工具应用向信息素养培育的范式迁移，着重培养学习者的数字化学习力与创新性思维 \cite{2018_教育部_教育部} 。《中国教育现代化 2035》进一步强调，信息化人才应具备技术整合能力、跨学科问题解决能力及人机协同创新能力，形成 “数字素养+学科核心素养” 的复合型能力结构 \cite{2019_中共中央_国务院_中国教育现代化2035}。

因此，新时代人才培养的现实需求迫切要求教师将信息技术深度融入教育实践，通过开发数字化教学素材、革新教育实施方法、探索育人模式转型，形成信息赋能与育人方式创新的联动机制，以切实回应时代发展对新型人才培养提出的全新挑战。在此背景下，高中数学实验教学资源开发需着力培育学生数学建模能力与数字化创新能力，通过构建虚实融合的实验环境，发展基于实例实操的科学探究素养，切身实际地提升学生综合素质，也为广大教师群体提供教学借鉴。

\subsection{高中数学课程改革的需要}

《普通高中数学课程标准（2017 年版 2020 年修订）》明确提出 “强化信息技术与数学课程的深度融合，推动数学实验教学资源创新开发”，强调通过数字化工具促进数学抽象概念的具象化与动态化呈现，用数字化工具把抽象数学概念转化成看得见、摸得着的动态形式，引导学生主动探究 \cite{2020_中华人民共和国教育部_普通高中数学课程标准}。

传统教学模式下的数学实验，作为学生理解数学概念、掌握数学方法、培养逻辑推理能力的直观途径，其动手操作与直观感知的价值不容小觑。然而，当前传统教学模式下的数学实验面临着数据记录繁琐、分析效率低下、结果呈现不够直观等挑战，难以满足师生对于实验精准度、效率及可视化方面的更高追求。相比之下，数字化实验教学资源凭借其强大的数据处理能力、即时反馈机制、直观的数据可视化效果以及精准捕捉数学变化特征的优势，展现出传统教学模式难以比拟的高效性与便捷性。例如，利用 GeoGebra 软件进行的函数图象绘制、借助希沃白板展现圆锥曲线的形成等实验，不仅简化了实验步骤，提升了实验效率，还极大地丰富了实验的表现形式，使得抽象的数学概念得以生动展现，有助于学生深化理解并激发创新思维。

数字化实验作为一种先进的信息技术手段，能够构建出一种全新的数学学习环境，通过实验教学真正落实学生创新思维与实践操作能力的同步提升，实现了从知识传授向能力培养的根本转变。实验教学资源的数字化转型不仅是教育现代化的必然趋势，更是促进学生全面发展、提升数学教育质量的关键路径。


\section{研究综述}

\subsection{国内外研究现状}

国际研究动态显示，技术工具与数学实验的融合已成为该领域的重要发展方向。现有文献多以信息技术实验资源在数学教学中的应用与分析研究为主，少有数学实验教学资源开发类文献。以下重在罗列阶段代表性研究成果。

20 世纪 80 年代中期，伴随计算机及软件技术的革命性进步，美国教育界发起了覆盖全国的深度微积分课程改革。此次改革以 “精简而直观的微积分教学” 为核心理念，其创新路径在于将数学实验教学模式系统化引入课堂教学实践\cite{2004_陈晓瑛_国内外数学实验教学的现状分析与展望}。

二十一世纪初，Markus Hohenwarter 团队开发的 GeoGebra 系统通过动态几何环境重构了数学认知路径，这是 GeoGebra 系统第一次出现在大众视野\cite{2011_Velichova_Interactive_Maths_with_GeoGebra}。该教授于 2007 年提出，GeoGebra 作为一款免费简便的开源软件，可以作为资源开发载体应用于实际实验教学中，是专门为教学过程而设计地一款软件。随后不久，Scott Reid 经研究提出教师是在课堂上使用信息化资源的先行者，鼓励实验式交互教学 \cite{2002_Reid_The_integration_of_information_and_communication_technology_into_classroomteaching}。

近年来，Tola Bekene Bedada 和 M.F. Machaba 利用 GeoGebra 软件开发数学学习的循环模型，将现代化信息技术融合于数学课堂中，开展数学实验学习模式\cite{2024_Bedada_The_development_of_the_cycle_model_and_its_effect_on_mathematics_learning_using_GeoGebra_mathematical_software}。

国内进展方面，借鉴学习国际信息技术资源化教学模式，针对于实验资源开发与应用分析，诸多学者皆开展相关研究。

2018 年，张志勇在研究中指出，数学实验可视化的核心目标在于通过图形化或动态化的表征形式 （如几何图形、数学模型动画等），将数学领域抽象的知识要素转化为具象化认知载体，从而帮助学生建立对数学知识体系的直观、整体的认知架构 \cite{2018_张志勇_高中数学可视化教学}。

2020 年，黄亚昕通过整合 GeoGebra 的三维动态建模与空间可视化技术，针对立体几何教学的知识建构难点，构建了基于该软件功能特性的实验教学设计模型。研究以 “空间向量关系推导” “几何体截面动态生成” 等案例为载体，验证了 GeoGebra 在提升学生空间想象与几何论证能力方面的实践效能，凸显技术工具与学科实验的深度适配性\cite{2020_黄亚昕_基于GeoGebra的立体几何教学研究}。

2021 年，董林伟及其团队在研究与实践数学实验教学的基础上，提出 “做数学”的主张，即让学生在数学实验探究过程中主观能动地习得知识、掌握方法，进而发展思维\cite{2021_董林伟}。

2022 年，邓璇从教师层面出发，分析运用希沃白板辅助高中数学教学策略，剖析其演示、记录、交互、资源库四大功能在实际教学中的应用\cite{2022_邓璇_高中数学运用希沃白板辅助教学的策略}。

2023 年，李婧妍提出创设虚拟实验环境，深度融合信息技术与高中数学实验课程，模拟真实交互场景，以培养学生实验设计与数据分析能力 \cite{2023_李婧妍_谈信息技术与高中数学课程的整合}。

\subsection{研究现状总结}

随着信息科技新世纪的快速发展，信息技术与教学的融合已成为必然趋势。 国内外都在教学中借助信息科技开展实验教学设计，比较于传统教学模式皆有一定优效。且视国内高中数学教学而言，目前已有相关实验资源配套教学，但资源有一定局限性，一定程度上限制了教师发挥课堂主体的自由性，即要求教师按规定好的教学路线与流程开展教学。对于实验教学片段资源的设计与开发，作为片段教学供给教师作为课堂教学元件还有待补充，此份工作也艰巨而远大。

\section{研究目的及意义}

\subsection{研究目的}

本研究基于高中数学二维平面版块新授课一平面向量基本定理、三角函数图象平移伸缩变换与三维空间板块趣味课一圆锥曲线的形成的核心教学难点，深度融合动态数学软件 GeoGebra 与交互式教学平台希沃白板，致力于构建低成本、 高交互的数字化实验教学体系。针对传统教学中函数图象变换过程可视化不足、 空间截面认知抽象度高等现实困境，通过 GeoGebra 的参数化建模技术实现函数参数的实时调控与三维几何体的动态剖切的平面展示，使抽象的数学规律转化为可触摸的认知脚手架。同时，将希沃白板的双屏协同功能与 GeoGebra 软件的数据变换可视化功能创新性地应用于课堂互动，实现教师端理论推导与学生端实践操作的深度融合。

\subsection{研究意义}

\subsubsection{理论意义}

针对《普通高中数学课程标准（2017 年版 2020 年修订）》提出的“数学探究与建模”素养培育要求，本研究基于 GeoGebra 与希沃白板，突破传统数学实验资源交互性弱、生成性不足的桎梏。通过深入剖析平面向量基本定理、圆锥曲线生成机理等核心概念的认知困难，结合信息化发展的数字软件开发应用实验资源，为数学实验教学理论注入新技术时代内涵。

\subsubsection{实践意义}

研究开发的高中数学实验教学资源，显著降低动态数学实验的技术实施门槛。 通过 GeoGebra 资源制作将模型变换转化为可视化操作模块，使教师无需编程基础即可快速生成函数图象连续变换动画。在课堂教学层面，创新设计的 “圆锥曲线的形成” 实验提升对抽象几何概念的认知，较传统讲授模式具有显著性差异。 助力教师实现从 “知识传授者” 到 “认知搭建者” 的专业转型，同时让学生从动态数学实验中深刻体会到数学探究的乐趣，主观能动地建构数学知识框架，发现和揭示数学现象背后的本质规律，有利于培养学生的问题意识、数学建模能力和批判性思维，提升其数学核心素养。

\section{研究内容及方法}

\subsection{研究内容}

第一，系统分析我国高中数学实验教学现状，明确需要攻克的难题。研读近年来中共中央办公厅、教育部颁布的教育改革政策文件，及《普通高中数学课程标准》，明晰新时代对创新型数学人才培养的核心诉求。在查阅大量国内外数学实验教学资料后，对比分析传统实验与信息化软件支持的实验研究进展，明确我国数字化数学实验存在的不足之处，针对现存短板，提出引入低成本开源数学工具 GeoGebra、希沃白板等作为实验载体，构建普惠型数字化实验开发路径。

第二，对相关概念界定，依托建构主义学习理论与认知发展理论，为高中数学实验教学的资源开发、活动设计及实施策略提供理论支撑与方法指导。通过文献阅读、网络媒体与切身实际感受，了解高中数学教学现状，明确教师与学生的信息化诉求。通过师生分析与教学设计分析，确立信息化数字实验的继开发迫在眉睫。

第三，基于 GeoGebra 软件与希沃白板智能交互系统，选择性开发部分高中实验教学资源——“平面向量基本定理”、“三角函数图象平移伸缩变换”、“圆锥曲线的形成”，涉及课程新授讲解、课堂实操演示、课外兴趣拓展三大模块。开发过程严格遵循科学性、适切性、可视化等原则，致力于开发普适性、普惠性的实验教学资源。

第四，在已开发实验资源基础上，进行相应实验讲义设计与教学片段设计，并开展实际模拟教学，将理论付诸于实践。并在实际模拟教学实践中分析反思，发现开发实验的创新价值与不足之处。

最后，对本研究进行系统地总结，同时，对后续研究进行展望。

\subsection{研究方法}

\subsubsection{文献研究法}

通过搜集文献，梳理文件与分析，了解信息化数字软件在开发数学教学实验中的研究进展与现状。系统的总结现有研究开发成果的优势与不足，从而转化为本研究的切入点与创新点。

\subsubsection{实验探究法}

依据高中数学课程标准与内容模块，基于 GeoGebra 软件和希沃白板开发具体实验，形成关联式实验体系。并通过模拟课堂教学，预测实验教学的应用，进行相关分析。

\subsubsection{案例分析法}

以培养学生核心素养为导向，对开发的高中数学实验教学资源进行片段教学设计，实现理论与现实的融合，为投入实际应用架构桥梁，以促进学生素养的塑造与提升。
